% ============================================================
% CAPÍTULO 5 — Razones de No Sostenibilidad
% ============================================================
\chapter{Razones de No Sostenibilidad}
\label{cap:razones}

\section{Resumen estructural}

A partir del an\'alisis del cap\'itulo~\ref{cap:sostenibilidad}, se
identifican cuatro razones estructurales que motivan la no
sostenibilidad del proyecto en su forma actual. Estas razones son
\textbf{estructurales, no personales}, y comunes en proyectos
estudiantiles que escalan en complejidad t\'ecnica.

\subsection*{R1. Escala humana}

La complejidad multistack del sistema requiere un equipo
multidisciplinar de aproximadamente cuatro personas con dedicaci\'on
parcial (60--80\,\% FTE total distribuido). La rotaci\'on natural de
miembros en una asociaci\'on estudiantil --- con permanencia media de
2--4 a\~nos --- no permite mantener ese nivel de capacidad t\'ecnica
acumulada de forma sostenida sin un programa expl\'icito de transferencia
de conocimiento que la asociaci\'on no tiene actualmente capacidad de
implementar.

\subsection*{R2. Escala econ\'omica}

El coste operativo recurrente m\'inimo
($\sim$459\,\euro/a\~no) excede por completo la financiaci\'on
recurrente disponible (0\,\euro/a\~no). No se trata de un desfase
parcial: \textbf{no existe ninguna partida comprometida} para el gasto
corriente del proyecto. La \'unica financiaci\'on institucional recibida
--- los 900\,\euro\ del Proyecto Docente de la ETSIA --- fue una
dotaci\'on puntual de inversi\'on en material, no una asignaci\'on
recurrente. El \textit{gap} requiere, por tanto, crear desde cero una
fuente de ingresos sostenida que la asociaci\'on no puede asegurar bajo
el formato actual.

Un dato ilustra el problema mejor que cualquier proyecci\'on: el
\'ultimo \textit{commit} del repositorio es de marzo de 2026, pero la
factura de Google Cloud sigui\'o llegando cada mes durante cinco meses
m\'as, hasta que la cuenta se cerr\'o de forma expresa. \textbf{El coste
de un sistema desplegado no se detiene cuando se detiene el desarrollo};
solo se detiene cuando alguien lo apaga. Esa asimetr\'ia es precisamente
lo que convierte una infraestructura sin equipo en un pasivo, y explica
por qu\'e el cierre ordenado que propone este informe no es una
formalidad administrativa sino una medida con efecto econ\'omico
inmediato.

\subsection*{R3. Escala de complejidad}

El proyecto integra tecnolog\'ias de seis dominios distintos
(\textit{embedded}, \textit{networking}, backend, frontend, DevOps,
agronom\'ia). Cada dominio aisladamente es manejable, pero la
integraci\'on y mantenimiento simult\'aneo demandan especializaci\'on
multidisciplinar dif\'icil de mantener en el tiempo en un contexto
estudiantil donde los miembros priorizan, l\'ogicamente, sus estudios.

\subsection*{R4. Single Point of Failure}

El desarrollo se ha concentrado en un \'unico miembro, creando una
dependencia estructural incompatible con la rotaci\'on natural de la
asociaci\'on. La distribuci\'on real del conocimiento del sistema
implica que, sin este miembro, el resto de la asociaci\'on no tiene
capacidad operativa sobre el conjunto del sistema.

\section{\textquestiondown Por qu\'e cerrar ahora y no antes?}

Estas limitaciones eran parcialmente conocidas desde el inicio del
proyecto. La decisi\'on de cierre formal \textbf{ahora}, frente a una
continuaci\'on indefinida, responde a las siguientes consideraciones:

\begin{itemize}[leftmargin=*]
    \item \textbf{El sistema ha alcanzado un hito relevante}: la
    funcionalidad completa est\'a documentada y validada por capa.
    El cierre se realiza \textbf{con valor demostrable}, no en estado
    intermedio.

    \item \textbf{La identificaci\'on clara de las limitaciones permite
    cierre razonado}: no es abandono por agotamiento, sino conclusi\'on
    fundamentada en datos t\'ecnicos y econ\'omicos espec\'ificos.

    \item \textbf{Mantener el proyecto activo bajo condiciones no
    sostenibles generar\'ia coste creciente sin valor proporcional},
    y eventualmente da\~no reputacional al no poder cumplir
    expectativas operativas.

    \item \textbf{El cierre maduro preserva la relaci\'on institucional}
    y deja abiertas opciones de continuidad bajo nuevos formatos
    (caso de estudio docente, base para TFG/TFM, transferencia
    eventual a un grupo de investigaci\'on).
\end{itemize}

\begin{warningbox}
Reconocer cu\'ando un proyecto debe cerrar es una competencia
profesional reconocida en la literatura de gesti\'on de proyectos
(Edmondson, 2023). En el contexto de proyectos docentes, esta
capacidad de evaluaci\'on cr\'itica es uno de los aprendizajes m\'as
valiosos que la experiencia ha aportado a los miembros participantes.
\end{warningbox}
